\chapter{Representaties van eindige groepen}\label{ch:b3:representations}

Om een abstracte groep te begrijpen, laat je haar werken op een
vectorruimte en zet je lineaire algebra --- eigenwaarden, sporen,
inproducten --- op die \omterm{def:b3:groups:action}{werking} los. Dat programma, de
\emph{representatietheorie}, is verbluffend doeltreffend voor
eindige groepen over $\C$: elke \omterm{def:b3:representations:rep}{representatie} valt uiteen in
irreducibele (Maschke), de irreducibele liggen vast door hun
\emph{\omterm{def:b3:representations:character}{karakters}} (sporen), en die \omterm{def:b3:representations:character}{karakters} voldoen aan
\omterm{thm:b3:representations:orthogonality}{orthogonaliteitsrelaties} die het rekenwerk mechanisch maken. Dit
hoofdstuk bouwt dat rekentuig op met zijn eerste meesterstukken
--- de \omterm{def:b3:representations:table}{karaktertafels} van de kleine groepen --- en de
weekendopgave oogst een stelling die ver buiten het bereik van de
kale groepentheorie ligt: de $p^aq^b$-stelling van Burnside.
Overal is $G$ een eindige groep en zijn alle vectorruimten
eindigdimensionaal over $\C$.

\section{Representaties, Maschke, Schur}

\begin{definition}\label{def:b3:representations:rep}
Een \emph{representatie}\index{representatie} van $G$ is een
morfisme $\rho \colon G \to GL(V)$ voor een $\C$-vectorruimte $V$;
$\dim V$ heet de \emph{graad} ervan. Een deelruimte $W \subseteq
V$ heet \emph{invariant} als $\rho(g)W \subseteq W$ voor alle $g$;
beperking maakt van $W$ een \emph{deelrepresentatie}. $\rho$ heet
\emph{irreducibel}\index{irreducibele representatie} als $V \neq
0$ en haar enige invariante deelruimten $0$ en $V$ zijn. Een
\emph{morfisme} tussen $(\rho, V)$ en $(\sigma, W)$ is een
lineaire $f\colon V \to W$ met $f\rho(g) = \sigma(g)f$ voor alle
$g$ (equivariantie); bijectieve $f$ heten \emph{isomorfismen}.
\end{definition}

\begin{example}\label{ex:b3:representations:examples}
(a) Graad $1$: de morfismen $G \to \C^\times$.
(b) De \emph{reguliere representatie}\index{reguliere
representatie}: $V = \C^G$ met basis $(e_h)_{h \in G}$ en
$\rho(g)e_h = e_{gh}$; graad $\abs G$.
(c) Een permutatiewerking van $G$ op een eindige verzameling $X$
geeft de \emph{permutatierepresentatie} op $\C^X$: $\rho(g)e_x =
e_{g\cdot x}$.
(d) $S_n$ werkt op $\C^n$ door de coördinaten te permuteren; het
hypervlak $\{\sum x_i = 0\}$ is invariant: de
\emph{standaardrepresentatie}, van graad $n - 1$.
\end{example}

\begin{theorem}[Maschke]\label{thm:b3:representations:maschke}
Elke invariante deelruimte $W$ van een \omterm{def:b3:representations:rep}{representatie} $(V, \rho)$
bezit een invariant complement. Bijgevolg is elke \omterm{def:b3:representations:rep}{representatie}
een directe som van irreducibele
(\emph{halfenkelvoudigheid}).\index{stelling van Maschke}
\end{theorem}

\begin{proof}
Zij $p \colon V \to V$ een \emph{willekeurige} projectie met beeld
$W$ (kies een willekeurig complement). Middel haar over de groep
uit:
\[
\tilde p = \frac1{\abs G}\sum_{g \in G} \rho(g)\,p\,\rho(g)^{-1}.
\]
Elke term beeldt $V$ in $W$ af ($W$ is invariant) en houdt $W$
puntsgewijs vast: voor $w \in W$ is $\rho(g)^{-1}w \in W$, houdt
$p$ dat vast, en brengt $\rho(g)$ het terug --- dus is $\tilde p$
opnieuw een projectie op $W$. Ze is equivariant: voor $h \in G$
herindexeert $\rho(h)\tilde p\rho(h)^{-1}$ dezelfde som. Bijgevolg
is $\ker \tilde p$ een invariant complement van $W$. Herhaald
toepassen op de summanden (eindige dimensie) ontbindt $V$ in
irreducibele.
\end{proof}

\begin{theorem}[Lemma van Schur]\label{thm:b3:representations:schur}
Zij $f \colon V \to W$ een morfisme van \emph{irreducibele}
\omterm{def:b3:representations:rep}{representaties}. Dan is $f = 0$ of is $f$ een isomorfisme; en als
$(V,\rho) = (W,\sigma)$, dan is $f = \lambda\,\mathrm{id}$ voor
zekere $\lambda \in \C$. Bijgevolg is $\dim\operatorname{Hom}_G(V,
W)$ gelijk aan $1$ als $V \cong W$, en anders $0$.\index{lemma van
Schur}
\end{theorem}

\begin{proof}
$\ker f$ en $\operatorname{im} f$ zijn invariant (equivariantie),
dus elk van beide is $0$ of het geheel: ofwel $f = 0$, ofwel is
$f$ injectief met volledig beeld. Is $V = W$, dan heeft $f$ een
eigenwaarde $\lambda$ ($\C$ is \omterm{def:b3:galois:closure}{algebraïsch gesloten}); $f -
\lambda\,\mathrm{id}$ is een niet-injectief morfisme $V \to V$ en
dus $0$. Voor de dimensietelling wanneer $V \cong W$: leg één
isomorfisme $u$ vast; elk morfisme $f$ geeft het endomorfisme
$u^{-1}f = \lambda\,\mathrm{id}$, dus $f = \lambda u$.
\end{proof}

\section{Karakters en orthogonaliteit}

\begin{definition}\label{def:b3:representations:character}
Het \emph{karakter}\index{karakter} van $(V, \rho)$ is
$\chi_\rho(g) = \operatorname{tr}\rho(g)$. Er geldt $\chi_\rho(e)
= \dim V$ en $\chi_\rho(hgh^{-1}) = \chi_\rho(g)$ (sporen zijn
invariant onder conjugatie): karakters zijn dus
\emph{klassefuncties}\index{klassefunctie} --- elementen van de
ruimte $\mathcal{CF}(G)$ van functies die constant zijn op de
\omterm{ex:b3:groups:actions}{conjugatieklassen}, voorzien van het hermitische inproduct
\[
\langle \varphi, \psi\rangle = \frac1{\abs G}\sum_{g \in G}
\overline{\varphi(g)}\,\psi(g).
\]
\end{definition}

\begin{proposition}\label{prop:b3:representations:charbasics}
$\rho(g)$ is diagonaliseerbaar met eenheidswortels als
eigenwaarden; $\chi_\rho(g^{-1}) = \overline{\chi_\rho(g)}$, en
$\abs{\chi_\rho(g)} \leq \chi_\rho(e)$ met gelijkheid precies
wanneer $\rho(g)$ scalair is. \omterm{def:b3:representations:character}{Karakters} tellen op bij directe
sommen: $\chi_{V \oplus W} = \chi_V + \chi_W$.
\end{proposition}

\begin{proof}
Er geldt $\rho(g)^N = \mathrm{id}$ voor $N = \abs G$ (Lagrange):
$\rho(g)$ annihileert $X^N - 1$, die uiteenvalt met enkelvoudige
wortels, zodat $\rho(g)$ diagonaliseerbaar is met eigenwaarden
$\lambda_i \in \mu_N$. Dan is $\chi(g^{-1}) = \sum \lambda_i^{-1}
= \sum\bar\lambda_i = \overline{\chi(g)}$, en $\abs{\chi(g)} =
\abs{\sum\lambda_i} \leq \dim V$, met gelijkheid in de
driehoeksongelijkheid precies wanneer alle $\lambda_i$ gelijk
zijn, dat wil zeggen $\rho(g) = \lambda\,\mathrm{id}$. De
additiviteit bij directe sommen: sporen per blok.
\end{proof}

\begin{lemma}\label{lem:b3:representations:homtrace}
Zij $(V, \rho)$ en $(W, \sigma)$ \omterm{def:b3:representations:rep}{representaties}. De
middelingsoperator op $\operatorname{Hom}(W, V)$,
\[
c(A) = \frac1{\abs G}\sum_{g}\rho(g)\,A\,\sigma(g)^{-1},
\]
is een projectie op $\operatorname{Hom}_G(W, V)$, en voor de
afbeelding $\Phi_g \colon A \mapsto \rho(g)A\sigma(g)^{-1}$ geldt
$\operatorname{tr}\Phi_g =
\chi_\rho(g)\,\overline{\chi_\sigma(g)}$.
\end{lemma}

\begin{proof}
$c(A)$ is equivariant (herindexeer de som als bij Maschke), en $c$
houdt equivariante afbeeldingen vast (elke term is gelijk aan
$A$): dus is $c$ een projectie met beeld $\operatorname{Hom}_G(W,
V)$. Spoor: in bases is $\Phi_g(A) = BAC$ met $B = \rho(g)$ en $C
= \sigma(g)^{-1}$; op de basis $(E_{kl})$ van matrices is
$BE_{kl}C = \sum_{m,n} b_{mk}c_{ln}E_{mn}$, zodat de coëfficiënt
van $E_{kl}$ in $\Phi_g(E_{kl})$ gelijk is aan $b_{kk}c_{ll}$:
$\operatorname{tr}\Phi_g = \sum_{k,l}b_{kk}c_{ll} =
\operatorname{tr}(B) \operatorname{tr}(C) =
\chi_\rho(g)\chi_\sigma(g^{-1})$, en $\chi_\sigma(g^{-1}) =
\overline{\chi_\sigma(g)}$.
\end{proof}

\begin{theorem}[Eerste orthogonaliteitsrelaties]
\label{thm:b3:representations:orthogonality}
Zij $\rho$ en $\sigma$ \emph{\omterm{def:b3:representations:rep}{irreducibel}}. Dan
is\index{orthogonaliteit van karakters}
\[
\langle\chi_\sigma, \chi_\rho\rangle =
\begin{cases}
1 & \text{als } \rho \cong \sigma,\\
0 & \text{anders:}
\end{cases}
\]
de irreducibele \omterm{def:b3:representations:character}{karakters} vormen een orthonormale familie in
$\mathcal{CF}(G)$.
\end{theorem}

\begin{proof}
Het spoor van een projectie is de dimensie van haar beeld:
\[
\dim\operatorname{Hom}_G(W, V) = \operatorname{tr} c
= \frac1{\abs G}\sum_g \operatorname{tr}\Phi_g
= \frac1{\abs G}\sum_g \chi_\rho(g)\overline{\chi_\sigma(g)}
= \langle \chi_\sigma, \chi_\rho\rangle,
\]
en het lemma van Schur evalueert het linkerlid tot $\delta_{\rho
\cong \sigma}$.
\end{proof}

\begin{corollary}\label{cor:b3:representations:multiplicity}
Ontbind $V \cong \bigoplus_i V_i^{\oplus m_i}$ in verschillende
irreducibele ($V_i \not\cong V_j$). Dan is $m_i = \langle
\chi_{V_i}, \chi_V\rangle$: de multipliciteiten --- en dus de
\omterm{def:b3:representations:rep}{representatie} op isomorfie na --- liggen door het \omterm{def:b3:representations:character}{karakter} vast.
Bovendien is $\langle \chi_V, \chi_V\rangle = \sum_i m_i^2$; in
het bijzonder is $V$ \omterm{def:b3:representations:rep}{irreducibel} dan en slechts dan als
$\langle\chi_V, \chi_V\rangle = 1$.
\end{corollary}

\begin{proof}
Er geldt $\chi_V = \sum m_i\chi_{V_i}$
(\cref{prop:b3:representations:charbasics}); neem inproducten met
elke $\chi_{V_i}$ en gebruik de orthonormaliteit. Twee
\omterm{def:b3:representations:rep}{representaties} met gelijke \omterm{def:b3:representations:character}{karakters} hebben gelijke
multipliciteiten en zijn dus isomorf.
\end{proof}

\begin{theorem}[De reguliere representatie]
\label{thm:b3:representations:regular}
Zij $\chi_1, \dots, \chi_r$ de verschillende irreducibele
\omterm{def:b3:representations:character}{karakters}, van graden $n_i = \chi_i(e)$. De \omterm{ex:b3:representations:examples}{reguliere
representatie} ontbindt met multipliciteiten $m_i = n_i$; bijgevolg
is
\[
\sum_{i=1}^{r} n_i^2 = \abs G,
\qquad
\sum_i n_i\chi_i(g) = 0 \quad (g \neq e).
\]
\end{theorem}

\begin{proof}
Het reguliere \omterm{def:b3:representations:character}{karakter} is $\chi_{\mathrm{reg}}(g) = \#\{h : gh =
h\}$, dus $\abs G$ voor $g = e$ en $0$ anders. Bijgevolg is $m_i =
\langle \chi_i,\chi_{\mathrm{reg}}\rangle = \frac1{\abs
G}\overline{\chi_i(e)}\,\abs G = n_i$. Evalueren van
$\chi_{\mathrm{reg}} = \sum n_i\chi_i$ in $e$ en in $g \neq e$
geeft de twee getoonde identiteiten.
\end{proof}

\begin{theorem}\label{thm:b3:representations:numberirr}
De irreducibele \omterm{def:b3:representations:character}{karakters} vormen een orthonormale \emph{basis} van
$\mathcal{CF}(G)$: het aantal \omterm{def:b3:representations:rep}{irreducibele representaties} is
gelijk aan het aantal \omterm{ex:b3:groups:actions}{conjugatieklassen} van $G$.
\end{theorem}

\begin{proof}
Alleen de volledigheid rest: zij $f \in \mathcal{CF}(G)$
orthogonaal met elke $\chi_i$; we tonen aan dat $f = 0$. Stel voor
een \omterm{def:b3:representations:rep}{representatie} $(V,\rho)$ de operator $T_{f,\rho} =
\frac{1}{\abs G} \sum_g \overline{f(g)}\,\rho(g)$. Die is
equivariant: voor $h \in G$ is
\[
\rho(h)T_{f,\rho}\rho(h)^{-1} = \frac1{\abs G}\sum_g
\overline{f(g)}\rho(hgh^{-1})
= \frac1{\abs G}\sum_{g'}\overline{f(h^{-1}g'h)}\rho(g') =
T_{f,\rho}
\]
($f$ is een \omterm{def:b3:representations:character}{klassefunctie}). Is $\rho$ \omterm{def:b3:representations:rep}{irreducibel} van graad $n$,
dan geeft Schur dat $T_{f,\rho} = \lambda\,\mathrm{id}$ met
\[
\lambda = \frac{\operatorname{tr}T_{f,\rho}}{n}
= \frac{1}{n\abs G}\sum_g \overline{f(g)}\,\chi_\rho(g)
= \frac1n\,\langle f, \chi_\rho\rangle = 0 .
\]
Dus is $T_{f,\rho} = 0$ op elke \omterm{def:b3:representations:rep}{irreducibele representatie}, en
daarmee (directe sommen) op \emph{elke} \omterm{def:b3:representations:rep}{representatie} --- in het
bijzonder op de reguliere. Pas hem toe op de basisvector $e_e$: $0
= T_{f,\mathrm{reg}}e_e = \frac1{\abs
G}\sum_g\overline{f(g)}e_g$, wat elke $\overline{f(g)} = 0$
afdwingt. De orthonormale familie $(\chi_i)$ spant dus
$\mathcal{CF}(G)$ op, en de dimensie daarvan is het aantal
\omterm{ex:b3:groups:actions}{conjugatieklassen}.
\end{proof}

\begin{corollary}[Kolomorthogonaliteit]
\label{cor:b3:representations:column}
Voor $g, h \in G$ geldt
\[
\sum_{i=1}^{r}\overline{\chi_i(g)}\,\chi_i(h) =
\begin{cases}
\abs{Z_G(g)} & \text{als } g \text{ en } h \text{ geconjugeerd
zijn},\\
0 & \text{anders.}
\end{cases}
\]
\end{corollary}

\begin{proof}
Zij $g_1, \dots, g_r$ representanten van de klassen en $c_j$ de
klassegroottes. De $r \times r$-matrix $U_{ij} =
\sqrt{c_j/\abs G}\;\chi_i(g_j)$ heeft orthonormale rijen
(\cref{thm:b3:representations:orthogonality} per klasse
geschreven: $\sum_j \frac{c_j}{\abs
G}\chi_i(g_j)\overline{\chi_{i'}(g_j)} = \delta_{ii'}$), dat wil
zeggen $UU^* = I$; en een vierkante matrix met $UU^* = I$ voldoet
ook aan $U^*U = I$: de kolommen zijn orthonormaal, wat zich
ontvouwt tot de getoonde identiteit ($\abs G/c_j =
\abs{Z_G(g_j)}$, de betrekking tussen \omterm{def:b3:groups:action}{baan} en \omterm{def:b3:groups:action}{stabilisator} voor de
conjugatie).
\end{proof}

\begin{proposition}[Eendimensionale karakters; optillen]
\label{prop:b3:representations:onedim}
(a) $G$ is abels dan en slechts dan als al haar \omterm{def:b3:representations:rep}{irreducibele
representaties} graad $1$ hebben; en voor elke $G$ is het aantal
\omterm{def:b3:representations:character}{karakters} van graad $1$ gelijk aan $[G : D(G)]$ (het zijn de
\omterm{def:b3:representations:character}{karakters} van de abelianisering).
(b) Is $N \trianglelefteq G$, dan tillen de \omterm{def:b3:representations:rep}{irreducibele
representaties} van $G/N$ zich (door samenstelling met $G \to G/N$)
op tot precies die \omterm{def:b3:representations:rep}{irreducibele representaties} van $G$ waarvan de
kern $N$ bevat.
\end{proposition}

\begin{proof}
(a) Is $G$ abels, dan is elke klasse een singleton: $r = \abs G$,
en $\sum n_i^2 = \abs G$ dwingt alle $n_i = 1$ af; zijn omgekeerd
alle $n_i = 1$, dan is de \omterm{ex:b3:representations:examples}{reguliere representatie} een som van
eendimensionale, zodat $\rho_{\mathrm{reg}}(G)$ simultaan
diagonaliseerbaar en dus commutatief is, en $\rho_{\mathrm{reg}}$
is trouw: $G$ is abels. \omterm{def:b3:representations:rep}{Representaties} van graad $1$ zijn
morfismen $G \to \C^\times$ met abels doel: ze factoriseren over
$G^{\mathrm{ab}} = G/D(G)$ (\cref{exo:b3:groups:9}), en het aantal
verschillende \omterm{def:b3:representations:character}{karakters} van de abelse groep $G^{\mathrm{ab}}$ is
$\abs{G^{\mathrm{ab}}}$ (zoveel klassen, alle van graad $1$). (b)
Samenstelling met de projectie bewaart de irreducibiliteit (de
invariante deelruimten corresponderen), en een \omterm{def:b3:representations:rep}{representatie} die
triviaal is op $N$ factoriseert over het quotiënt
(\cref{thm:b3:groups:firstiso}).
\end{proof}

\section{Karaktertafels}

\begin{definition}\label{def:b3:representations:table}
De \emph{karaktertafel}\index{karaktertafel} van $G$ is de $r
\times r$-matrix $\bigl(\chi_i(g_j)\bigr)$: de rijen zijn
geïndexeerd door de irreducibele \omterm{def:b3:representations:character}{karakters}, de kolommen door de
\omterm{ex:b3:groups:actions}{conjugatieklassen} (met hun groottes erbij vermeld). De rijen zijn
orthonormaal voor het gewogen product en de kolommen zijn
orthogonaal (\cref{cor:b3:representations:column}): de tafel is
zwaar overbepaald, en juist daarom berekenbaar.
\end{definition}

\begin{example}[De tafel van $S_3$]\label{ex:b3:representations:s3}
Klassen: $e$ (grootte 1), de transposities (3), de $3$-cykels (2);
dus $r = 3$ \omterm{def:b3:representations:rep}{irreducibele representaties}, van graden $n_i$ met
$\sum n_i^2 = 6$: $1, 1, 2$. Graad $1$: het triviale \omterm{def:b3:representations:character}{karakter}
$\mathbf 1$ en de signatuur $\varepsilon$. De laatste rij volgt
uit de \omterm{thm:b3:representations:orthogonality}{kolomorthogonaliteit} (of uit $\chi_{\mathrm{std}} =
\chi_{\mathrm{perm}} - \mathbf 1$):
\[
\begin{array}{c|ccc}
S_3 & e & (1\,2)\ [3] & (1\,2\,3)\ [2]\\
\hline
\mathbf 1 & 1 & 1 & 1\\
\varepsilon & 1 & -1 & 1\\
\chi_{\mathrm{std}} & 2 & 0 & -1
\end{array}
\]
Controle: $\langle\chi_{\mathrm{std}},\chi_{\mathrm{std}}\rangle =
\frac{1}{6}(4 + 0 + 2) = 1$: \omterm{def:b3:representations:rep}{irreducibel}.
\end{example}

\begin{example}[De tafel van $S_4$]\label{ex:b3:representations:s4}
Klassen: $e$ [1], de transposities [6], de dubbele transposities
[3], de $3$-cykels [8] en de $4$-cykels [6]: vijf \omterm{def:b3:representations:rep}{irreducibele
representaties} met $\sum n_i^2 = 24$ en twee van graad $1$
($\mathbf 1, \varepsilon$; want $[S_4 : D(S_4)] = [S_4 : A_4] =
2$): graden $1, 1, 2, 3, 3$. Die van graad $2$ tilt zich op uit
$S_4/V \cong S_3$ (\cref{prop:b3:representations:onedim}(b), met
$V$ de viergroep van Klein); graad $3$: de standaardrepresentatie
en haar twist met $\varepsilon$:
\[
\begin{array}{c|ccccc}
S_4 & e\,[1] & (1\,2)\,[6] & (1\,2)(3\,4)\,[3] &
(1\,2\,3)\,[8] & (1\,2\,3\,4)\,[6]\\
\hline
\mathbf 1 & 1 & 1 & 1 & 1 & 1\\
\varepsilon & 1 & -1 & 1 & 1 & -1\\
\chi_2 & 2 & 0 & 2 & -1 & 0\\
\chi_{\mathrm{std}} & 3 & 1 & -1 & 0 & -1\\
\varepsilon\chi_{\mathrm{std}} & 3 & -1 & -1 & 0 & 1
\end{array}
\]
($\chi_{\mathrm{std}}(g) = \operatorname{fix}(g) - 1$; en
$\chi_2$ evalueert de $S_3$-tafel op het beeld van elke klasse
modulo $V$.) Alle rij- en kolomcontroles slagen --- twee ervan
uitvoeren is de opwarmer van \cref{exo:b3:representations:3}.
\end{example}

\begin{method}\label{met:b3:representations:table}
Om een \omterm{def:b3:representations:table}{karaktertafel} op te bouwen: (1) geef de \omterm{ex:b3:groups:actions}{conjugatieklassen}
met hun groottes; (2) tel de \omterm{def:b3:representations:character}{karakters} van graad $1$ via $G/D(G)$
en schrijf ze op; (3) haal de overige graden uit $\sum n_i^2 =
\abs G$ (combinatoriek met kleine gehele getallen); (4) vind
goedkope irreducibele \omterm{def:b3:representations:character}{karakters}: til op uit quotiënten, trek
$\mathbf 1$ af van permutatiekarakters (controleer
$\langle\chi,\chi\rangle = 1$), vermenigvuldig bekende \omterm{def:b3:representations:character}{karakters}
met die van graad $1$; (5) maak de onbekende rijen af met
\omterm{thm:b3:representations:orthogonality}{kolomorthogonaliteit} --- elke kolom staat loodrecht op de reeds
volledige kolommen, en de kolom van $e$ draagt de graden.
Controleer alles met een volledige orthogonaliteitsronde.
\end{method}

\begin{omfigure}
\begin{tikzpicture}[scale=0.52]
  \draw[thick] (0,0) rectangle (6,6);
  \fill[omDef!25] (0,5) rectangle (1,6);
  \fill[omProp!25] (1,4) rectangle (2,5);
  \fill[omThm!20] (2,0) rectangle (6,4);
  \draw (0,5) rectangle (1,6);
  \draw (1,4) rectangle (2,5);
  \draw (2,0) rectangle (6,4);
  \node[font=\small] at (0.5,5.5) {$1$};
  \node[font=\small] at (1.5,4.5) {$1$};
  \node[font=\small] at (4,2) {$M_2(\C)$};
\end{tikzpicture}

{\small De \omterm{ex:b3:representations:examples}{reguliere representatie} van $S_3$, blokdiagonaal
gemaakt: $\C[S_3] \cong \C \times \C \times M_2(\C)$, met
dimensies $1 + 1 + 4 = 6 = \abs{S_3}$. In het algemeen is $\C[G]
\cong \prod_i M_{n_i}(\C)$: de identiteit $\sum n_i^2 = \abs G$ is
een uitspraak over matrixblokken.}
\end{omfigure}

\section{Oefeningen}

\begin{exercise}[$\star$]\label{exo:b3:representations:1}
(a) Toon aan dat de irreducibele \omterm{def:b3:representations:character}{karakters} van $\Z/n\Z$ de
$\chi_k(\bar m) = \eu^{2\iu\pi km/n}$ met $k = 0, \dots, n-1$
zijn, en schrijf de \omterm{def:b3:representations:table}{karaktertafel} van $\Z/4\Z$ op.
(b) Ga beide \omterm{thm:b3:representations:orthogonality}{orthogonaliteitsrelaties} erop na --- en herken de
matrix: waar is dit boek haar eerder tegengekomen?
\end{exercise}

\begin{exercise}[$\star$]\label{exo:b3:representations:2}
Zij $G$ werkend op een eindige verzameling $X$ en $\chi$ het
\omterm{def:b3:representations:character}{karakter} van de permutatierepresentatie $\C^X$.
(a) Toon aan dat $\chi(g) = \abs{\operatorname{Fix}_X(g)}$ en
$\langle \mathbf 1, \chi\rangle = \#\{\text{banen}\}$ --- de
teltelling van Burnside (\cref{exo:b3:groups:5}) is een
karakterberekening.
(b) Stel dat de \omterm{def:b3:groups:action}{werking} transitief is, zodat $\chi = \mathbf 1 +
\psi$. Toon aan dat $\psi$ \omterm{def:b3:representations:rep}{irreducibel} is dan en slechts dan als
de \omterm{def:b3:groups:action}{werking} \emph{$2$-transitief} is (transitief op geordende paren
verschillende punten). \emph{(Bereken $\langle\chi,\chi\rangle$
als het aantal \omterm{def:b3:groups:action}{banen} op $X \times X$.)}
(c) Besluit dat de standaardrepresentatie van $S_n$ ($n \geq 2$)
\omterm{def:b3:representations:rep}{irreducibel} is.
\end{exercise}

\begin{exercise}[$\star$]\label{exo:b3:representations:3}
Bouw de tafel van $S_3$ van nul af opnieuw op volgens
\cref{met:b3:representations:table}, en ga daarna twee rij- en
twee kolomorthogonaliteitsrelaties na in de $S_4$-tafel van
\cref{ex:b3:representations:s4}. Ontbind het permutatiekarakter
van $S_4$ werkend op $\{1,2,3,4\}$ en het \omterm{def:b3:representations:character}{karakter}
$\chi_{\mathrm{std}}^2$ (het puntsgewijze kwadraat) in
irreducibele \omterm{def:b3:representations:character}{karakters}.
\end{exercise}

\begin{exercise}[$\star\star$]\label{exo:b3:representations:4}
Bereken de \omterm{def:b3:representations:table}{karaktertafels} van $D_4$ en van $Q_8$. Besluit dat twee
niet-isomorfe groepen dezelfde \omterm{def:b3:representations:table}{karaktertafel} kunnen hebben ---
welke groepentheoretische gegevens legt de tafel bij dit paar toch
vast (de orde van het \omterm{ex:b3:groups:actions}{centrum}, de abelianisering, het aantal
involuties)? En welke daarvan legt ze juist \emph{niet} vast?
\end{exercise}

\begin{exercise}[$\star\star$]\label{exo:b3:representations:5}
\omterm{def:b3:representations:table}{Karaktertafel} van $A_4$: klassen $e$ [1], de dubbele transposities
[3], en \emph{twee} klassen van $3$-cykels [4] en [4].
(a) Verklaar de splitsing van de $3$-cykels (vergelijk de
\omterm{ex:b3:groups:actions}{centralisatoren} in $S_4$ en in $A_4$, als in
\cref{exo:b3:groups:11}).
(b) Bepaal de drie \omterm{def:b3:representations:character}{karakters} van graad $1$ (via $A_4/V \cong
\Z/3\Z$) en het \omterm{def:b3:representations:character}{karakter} van graad $3$ (beperk
$\chi_{\mathrm{std}}$ van $S_4$), en stel de tafel samen.
(c) Lees de \omterm{def:b3:groups:normal}{normaaldelers} van $A_4$ uit de tafel af (de kernen
$\{g : \chi(g) = \chi(e)\}$ en hun doorsneden).
\end{exercise}

\begin{exercise}[$\star\star$]\label{exo:b3:representations:6}
(a) Toon aan dat $\ker\chi = \{g : \chi(g) = \chi(e)\}$ de kern is
van de onderliggende \omterm{def:b3:representations:rep}{representatie}
(\cref{prop:b3:representations:charbasics}, het gelijkheidsgeval).
(b) Toon aan dat elke \omterm{def:b3:groups:normal}{normaaldeler} van $G$ een doorsnede is van
kernen van irreducibele \omterm{def:b3:representations:character}{karakters}. \emph{(Representeer $G/N$
trouw: haar \omterm{ex:b3:representations:examples}{reguliere representatie}.)}
(c) Leid af: $G$ is enkelvoudig dan en slechts dan als $\ker\chi_i
= \{e\}$ voor elk niet-triviaal \omterm{def:b3:representations:rep}{irreducibel} $\chi_i$ --- de
enkelvoudigheid is uit de \omterm{def:b3:representations:table}{karaktertafel} af te lezen.
\end{exercise}

\begin{exercise}[$\star\star$]\label{exo:b3:representations:7}
Graden bij $\abs G = 8$: toon aan dat een niet-abelse groep van
orde $8$ het gradenpatroon $(1,1,1,1,2)$ heeft en dat haar
\omterm{def:b3:representations:rep}{representatie} van graad $2$ trouw is. Toon algemener aan dat een
niet-abelse groep van orde $p^3$ het patroon $(1^{\,p^2}, p,
\dots, p)$ heeft, met $p^2$ enen en $p - 1$ \omterm{def:b3:representations:character}{karakters} van graad
$p$. \emph{(Gebruik $[G : D(G)]$ en $\sum n_i^2 = \abs G$; hier
heeft $D(G) = Z(G)$ orde $p$.)}
\end{exercise}

\begin{exercise}[$\star\star\star$]\label{exo:b3:representations:8}
Toon voor eindige groepen $G, H$ aan dat de \omterm{def:b3:representations:character}{klassefuncties}
$\chi(g)\psi(h)$ op $G \times H$, met $\chi, \psi$ irreducibele
\omterm{def:b3:representations:character}{karakters} van $G$ respectievelijk $H$, precies de irreducibele
\omterm{def:b3:representations:character}{karakters} van $G \times H$ zijn. \emph{(De orthonormaliteit is een
rechtstreekse berekening; de volledigheid volgt uit het tellen van
de klassen.)} Leid de \omterm{def:b3:representations:table}{karaktertafel} van $\Z/2\Z \times \Z/2\Z$ af
en bewijs \cref{prop:b3:representations:onedim}(a) opnieuw voor
eindige abelse groepen via de structuurstelling.
\end{exercise}

\begin{exercise}[$\star\star\star$]\label{exo:b3:representations:9}
De \omterm{def:b3:representations:table}{karaktertafel} van $A_5$ (klassen van grootte $1, 15, 20, 12,
12$ uit \cref{exo:b3:groups:11}):
(a) Toon aan dat de graden $1, 3, 3, 4, 5$ zijn \emph{(de enige
oplossing van $\sum n_i^2 = 60$ met $n_1 = 1$ en, wegens
\cref{exo:b3:representations:6}(c) samen met de enkelvoudigheid,
geen andere $n_i = 1$)}.
(b) Construeer het \omterm{def:b3:representations:character}{karakter} van graad $4$ (de permutatiewerking op
$5$ punten) en dat van graad $5$ (de \omterm{def:b3:groups:action}{werking} op de zes
Sylow-$5$-deelgroepen geeft graad $6 = 1 + 5$; controleer de
irreducibiliteit), en maak de twee rijen van graad $3$ af met
\omterm{thm:b3:representations:orthogonality}{kolomorthogonaliteit}: op de klassen van $5$-cykels verschijnen de
waarden $\frac{1\pm\sqrt5}2$ van de gulden snede.
(c) Ga aan de afgemaakte tafel na dat $A_5$ enkelvoudig is
(\cref{exo:b3:representations:6}(c)).
\end{exercise}

\begin{exercise}[$\star\star$]\label{exo:b3:representations:10}
Zij $\rho$ een \omterm{def:b3:representations:rep}{irreducibele representatie} van graad $n$ en $z \in
Z(G)$. Toon aan dat $\rho(z) = \lambda_z\,\mathrm{id}$ met
$\lambda \colon Z(G) \to \C^\times$ een morfisme (het
\emph{centrale \omterm{def:b3:representations:character}{karakter}}), en leid af dat $\abs{\chi(z)} = n$ voor
centrale $z$. Toepassing: heeft $G$ een trouwe \omterm{def:b3:representations:rep}{irreducibele
representatie}, dan is $Z(G)$ cyclisch.
\end{exercise}

\begin{exercise}[$\star\star$]\label{exo:b3:representations:11}
(Isotypische projecties) Zij $(V, \rho)$ een \omterm{def:b3:representations:rep}{representatie} van $G$
en $\chi_i$ een \omterm{def:b3:representations:rep}{irreducibel} \omterm{def:b3:representations:character}{karakter} van graad $n_i$. Definieer
\[
p_i = \frac{n_i}{\abs G}\sum_{g\in G}
\overline{\chi_i(g)}\,\rho(g) \;\in\; \mathcal L(V).
\]
(a) Toon aan dat $p_i$ $G$-equivariant is, en bereken haar
beperking tot een irreducibele deelrepresentatie $W \subseteq V$
met \omterm{def:b3:representations:character}{karakter} $\chi_j$: die is $\delta_{ij}\, \mathrm{id}_W$
\emph{(Schur; neem sporen om de scalair te bepalen)}.
(b) Leid af dat $p_i$ een projectie is op de som $V_i$ van alle
irreducibele deelrepresentaties met \omterm{def:b3:representations:character}{karakter} $\chi_i$ (de
\emph{isotypische component}), dat $\sum_ip_i = \mathrm{id}_V$, en
dat de ontbinding $V = \bigoplus_iV_i$ kanoniek is --- anders dan
de fijnere splitsing van elke $V_i$ in irreducibele.
(c) Schrijf voor de \omterm{ex:b3:representations:examples}{reguliere representatie} van $S_3$ en het
tekenkarakter $\varepsilon$ de projectie $p_\varepsilon$
expliciet op als element van de groepsalgebra, en controleer met
de hand dat $p_\varepsilon^2 = p_\varepsilon$.
\end{exercise}

\begin{exercise}[$\star\star$]\label{exo:b3:representations:12}
(Een tafel lezen) Van een zekere groep $G$ van orde $24$ is de
\omterm{def:b3:representations:table}{karaktertafel} gedeeltelijk bekend: ze heeft $5$ klassen, van
grootte $1, 6, 8, 6, 3$, en graden $1, 1, 2, 3, 3$.
(a) Reconstrueer de volledige tafel: eerst de twee lineaire
\omterm{def:b3:representations:character}{karakters} (één triviaal; het andere neemt de waarde $-1$ precies
op de klassen van grootte $6$ en $6$), dan het \omterm{def:b3:representations:character}{karakter} van graad
$2$ met \omterm{thm:b3:representations:orthogonality}{kolomorthogonaliteit} tegen de kolom van de identiteit, en
daarna evenzo de twee \omterm{def:b3:representations:character}{karakters} van graad $3$ (waarvan er één
$\chi_2\chi_4$ is).
(b) Herken $G$ ($\cong S_4$: vergelijk de klassen met de
cykeltypen), en haal met \cref{exo:b3:representations:6} de
\omterm{def:b3:groups:normal}{normaaldelers} uit de tafel: de kernen van $\chi_2$ (index $2$:
$A_4$) en van $\chi_3$ (de viergroep $V_4$), en verder niets
behalve $\{e\}$ en $G$.
(c) Leg uit hoe de tafel laat zien dat $G/V_4 \cong S_3$
\emph{(welke \omterm{def:b3:representations:character}{karakters} factoriseren over het quotiënt?)}.
\end{exercise}

\section{Probleem: de \texorpdfstring{$p^aq^b$}{paqb}-stelling van
Burnside}

\begin{problem}[{Weekendopgave --- oplosbaarheid van groepen van
orde $p^aq^b$}]\label{pb:b3:representations:1}
Burnside bewees in 1904 dat elke groep waarvan de orde hoogstens
twee priemfactoren heeft, \omterm{def:b3:groups:derived}{oplosbaar} is --- een uitspraak over
abstracte groepen waarvan alle bekende bewijzen een halve eeuw
lang via de karaktertheorie liepen. Dit probleem bouwt het bewijs
volledig op en brengt daarbij \cref{ch:b3:groups}
(\omterm{def:b3:groups:derived}{oplosbaarheid}), \cref{ch:b3:modules} (eindig voortgebrachte
$\Z$-modulen) en dit hoofdstuk samen. Overal zijn $\chi_1, \dots,
\chi_r$ de irreducibele \omterm{def:b3:representations:character}{karakters} van $G$ en $n_i = \chi_i(e)$.

\medskip
\noindent\textbf{Deel I --- Gehele \omterm{def:b3:galois:algebraic}{algebraïsche} getallen.} Een
\emph{geheel \omterm{def:b3:galois:algebraic}{algebraïsch} getal} is een wortel van een
\emph{monische} veelterm van $\Z[X]$.
\begin{enumerate}
  \item Toon aan dat $\alpha$ een geheel \omterm{def:b3:galois:algebraic}{algebraïsch} getal is dan
        en slechts dan als de ring $\Z[\alpha]$ een \omterm{def:b3:modules:free}{eindig
        voortgebracht} $\Z$-moduul is.
  \item Leid af dat de gehele \omterm{def:b3:galois:algebraic}{algebraïsche} getallen een deelring
        van $\C$ vormen. \emph{(Zijn $\Z[\alpha]$ en $\Z[\beta]$
        \omterm{def:b3:modules:free}{eindig voortgebracht}, dan ook $\Z[\alpha, \beta]$, en
        deelmodulen van eindig voortgebrachte $\Z$-modulen zijn
        \omterm{def:b3:modules:free}{eindig voortgebracht}, volgens
        \cref{thm:b3:modules:submodule} plus een
        presentatieargument --- of rechtstreeks: een deelmoduul
        van $\Z^n$ is vrij van rang $\leq n$.)}
  \item Toon aan dat een \emph{rationaal} geheel \omterm{def:b3:galois:algebraic}{algebraïsch} getal
        een geheel getal is. \emph{(De rationalewortelstelling.)}
  \item Toon aan dat elke karakterwaarde $\chi(g)$ een geheel
        \omterm{def:b3:galois:algebraic}{algebraïsch} getal is.
\end{enumerate}

\medskip
\noindent\textbf{Deel II --- De klassesomrelaties.} Leg een
irreducibele $(V, \rho)$ van graad $n$ met \omterm{def:b3:representations:character}{karakter} $\chi$ vast.
Stel voor een conjugatieklasse $C$: $S_C = \sum_{g \in C}\rho(g)
\in \mathcal L(V)$.
\begin{enumerate}[resume]
  \item Toon aan dat $S_C$ equivariant is, zodat $S_C =
        \omega_C\,\mathrm{id}$ met
        \[
        \omega_C = \frac{\abs C\,\chi(g_C)}{n}
        \qquad (g_C \in C \text{ een willekeurige representant}).
        \]
  \item Toon aan dat $S_CS_{C'} = \sum_{C''}
        a_{CC'C''}\,S_{C''}$, waarbij $a_{CC'C''} \in \N$ voor een
        vaste $z \in C''$ de paren $(x, y) \in C \times C'$ met
        $xy = z$ telt. Leid af dat de $\omega_C$ voldoen aan
        $\omega_C\,\omega_{C'} = \sum_{C''}
        a_{CC'C''}\,\omega_{C''}$.
  \item Besluit dat elke $\omega_C$ een geheel \omterm{def:b3:galois:algebraic}{algebraïsch} getal
        is. \emph{(Het $\Z$-moduul voortgebracht door $1$ en de
        $\omega_C$ is een eindig voortgebrachte ring; pas het
        criterium van vraag 1 toe --- preciezer: toon aan dat $M =
        \Z\text{-opspansel}(1, (\omega_C)_C)$ voldoet aan
        $\omega_{C_0} M \subseteq M$ en gebruik een truc met
        determinanten of Cayley--Hamilton, of het deelringargument
        van vraag 2.)}
  \item Leid de \emph{deelbaarheid van \omterm{thm:b3:galois:finitefields}{Frobenius}} af: $n_i$ deelt
        $\abs G$ voor elke irreducibele graad. \emph{(Bereken
        $\frac{\abs G}{n} = \frac{\abs G}{n}
        \langle\chi,\chi\rangle = \sum_C \omega_C\,
        \overline{\chi(g_C)}$: een geheel \omterm{def:b3:galois:algebraic}{algebraïsch} getal dat
        rationaal is.)}
\end{enumerate}

\medskip
\noindent\textbf{Deel III --- Het enkelvoudigheidscriterium van
Burnside.}
\begin{enumerate}[resume]
  \item Zij $\chi$ \omterm{def:b3:representations:rep}{irreducibel} van graad $n$ en $C$ een klasse met
        $\gcd(\abs C, n) = 1$. Toon met Bézout en de vragen 4--7
        aan dat $\frac{\chi(g_C)}{n}$ een geheel \omterm{def:b3:galois:algebraic}{algebraïsch} getal
        is.
  \item Stel bovendien $0 < \abs{\chi(g_C)} < n$. Toon aan dat dit
        onmogelijk is: het gehele \omterm{def:b3:galois:algebraic}{algebraïsche} getal $\alpha =
        \chi(g_C)/n$ heeft al haar \emph{toegevoegden} --- de
        getallen $\alpha_\sigma = \frac1n\sigma(\chi(g_C))$ voor
        $\sigma \in \operatorname{Gal}(\Q(\zeta_{\abs G})/\Q)$,
        elk een gemiddelde van $n$ eenheidswortels --- van modulus
        $\leq 1$, zodat het product $N = \prod_\sigma\alpha_\sigma$
        een rationaal geheel \omterm{def:b3:galois:algebraic}{algebraïsch} getal is met $0 < \abs N
        < 1$ --- verantwoord elke bewering, met beroep op
        \cref{thm:b3:galois:cyclotomicirred} voor de \omterm{def:b3:galois:galois}{Galoisgroep}
        en op het feit dat $\sigma$ eenheidswortels permuteert.
        Besluit: ofwel $\chi(g_C) = 0$, ofwel is $\rho(g_C)$
        scalair (\cref{prop:b3:representations:charbasics}).
  \item (Het criterium van Burnside) Zij $C \neq \{e\}$ een
        conjugatieklasse van \emph{priemmachtgrootte} $p^k > 1$,
        en stel dat $G$ enkelvoudig en niet-abels is.
        \omterm{thm:b3:representations:orthogonality}{Kolomorthogonaliteit} van de kolom van $C$ tegen die van
        $e$ geeft $1 + \sum_{i \geq 2} n_i\chi_i(g_C) = 0$. Toon
        aan dat een zeker niet-triviaal $\chi_i$ met $p \nmid n_i$
        voldoet aan $\chi_i(g_C) \neq 0$ \emph{(anders zou
        $\frac1p$ een geheel \omterm{def:b3:galois:algebraic}{algebraïsch} getal zijn)}; volgens
        vraag 10 is $\rho_i(g_C)$ dan scalair; leid daaruit een
        tegenspraak met de enkelvoudigheid af \emph{(de
        verzameling $g$ waarvoor $\rho_i(g)$ scalair is, is een
        \omterm{def:b3:groups:normal}{normaaldeler}; gebruik de trouwheid uit
        \cref{exo:b3:representations:6})}. Besluit: \emph{geen
        enkele enkelvoudige niet-abelse groep heeft een
        conjugatieklasse van priemmachtgrootte $> 1$.}
\end{enumerate}

\medskip
\noindent\textbf{Deel IV --- De stelling.}
\begin{enumerate}[resume]
  \item Zij $\abs G = p^aq^b$ met $a + b \geq 1$. Toon aan dat $G$,
        als ze enkelvoudig is, abels is: kies $z \neq e$ in het
        \omterm{ex:b3:groups:actions}{centrum} van een Sylow-$q$-deelgroep
        (\cref{thm:b3:groups:pfixed}) en beschouw de grootte van
        haar conjugatieklasse $[G : Z_G(z)]$, een macht van $p$
        (waarom?); pas vraag 11 toe.
  \item Besluit met inductie naar $\abs G$: \emph{elke groep van
        orde $p^aq^b$ is \omterm{def:b3:groups:derived}{oplosbaar}} (Burnside). Waarom breekt het
        argument bij drie priemgetallen --- en waarom moet het wel
        breken, gezien $\abs{A_5} = 2^2\cdot3\cdot5$?
\end{enumerate}

\medskip
\noindent\textbf{Deel V --- De \omterm{def:b3:representations:table}{karaktertafel} van $A_5$.} De
kleinste groep die de stelling van Burnside niet kan raken,
verdient haar volledige portret; alles hieronder gebruikt alleen
dit hoofdstuk plus \cref{exo:b3:groups:11}.
\begin{enumerate}[resume]
  \item Herinner uit \cref{exo:b3:groups:11} de vijf
        \omterm{ex:b3:groups:actions}{conjugatieklassen} van $A_5$: $\{e\}$, de $15$ dubbele
        transposities, de $20$ drietallige cykels, en twee klassen
        van elk $12$ vijftallige cykels, met representanten $c =
        (1\,2\,3\,4\,5)$ en $c^2$. Verklaar waarom de $5$-cykels
        in twee $A_5$-klassen uiteenvallen hoewel ze één enkele
        $S_5$-klasse vormen.
  \item Toon aan dat de irreducibele graden van $A_5$ precies $1,
        3, 3, 4, 5$ zijn: gebruik $\sum_in_i^2 = 60$ met $r = 5$
        klassen, en het feit dat $A_5$ \omterm{prop:b3:galois:perfect}{perfect} is ($D(A_5) =
        A_5$), zodat het triviale \omterm{def:b3:representations:character}{karakter} haar enige lineaire is;
        elimineer vervolgens elke andere multiverzameling
        \emph{(schrijf $59$ als som van vier kwadraten van gehele
        getallen $\geq 2$: ga na dat dat maar op één manier lukt
        met allemaal plausibele graden)}.
  \item Zij $\pi$ het permutatiekarakter van $A_5$ op $\{1, \dots,
        5\}$: $\pi(g) = \#\operatorname{Fix}(g)$, met waarden $5,
        1, 2, 0, 0$ op de vijf klassen. Bereken $\langle\pi,
        \mathbf 1\rangle$ en $\langle\pi, \pi\rangle$, en leid af
        dat $\chi_4 = \pi - \mathbf 1$ \omterm{def:b3:representations:rep}{irreducibel} is van graad
        $4$, met waarden $4, 0, 1, -1, -1$.
  \item Hetzelfde spel op de $10$ ongeordende paren $\{i, j\}$:
        daar zijn de aantallen vaste punten $10, 2, 1, 0, 0$.
        Bereken $\langle\pi_{10}, \pi_{10}\rangle$,
        $\langle\pi_{10}, \mathbf 1\rangle$ en $\langle\pi_{10},
        \chi_4\rangle$, leid de ontbinding $\pi_{10} = \mathbf 1 +
        \chi_4 + \chi_5$ af, en verkrijg het irreducibele
        $\chi_5$ van graad $5$ met waarden $5, 1, -1, 0, 0$.
  \item De twee resterende irreducibele \omterm{def:b3:representations:character}{karakters} $\chi_2$ en
        $\chi_3$ hebben graad $3$. \omterm{thm:b3:representations:orthogonality}{Kolomorthogonaliteit} (elke
        kolom $\neq e$ tegen de kolom van de identiteit, en elke
        kolom met zichzelf) legt hun waarden buiten de $5$-cykels
        vast: toon aan dat $\chi_2(g) = \chi_3(g) = -1$ op de
        dubbele transposities en $0$ op de drietallige cykels.
  \item Stel op de klassen van $5$-cykels $x = \chi_2(c)$ en $y =
        \chi_2(c^2)$; wegens de symmetrie mag men $\chi_3(c) = y$
        en $\chi_3(c^2) = x$ nemen. Leid uit de kolom van $c$
        gepaard met de kolom van de identiteit en met die van
        $c^2$ af dat $x + y = 1$ en $xy = -1$ (en controleer de
        waarde $x^2 + y^2 = 3$ die de kolom van $c$ met zichzelf
        geeft), waaruit
        \[
        \{x, y\} = \Bigl\{\frac{1 + \sqrt5}2,\ \frac{1 -
        \sqrt5}2\Bigr\} :
        \]
        de gulden snede en haar toegevoegde. Stel de volledige
        \omterm{def:b3:representations:table}{karaktertafel} van $A_5$ samen.
  \item Voer de controles uit: de rijnorm van $\chi_2$ is $1$
        (gebruik $\varphi^2 + \bar\varphi^2 = 3$),
        $\langle\chi_2, \chi_3\rangle = 0$, en de deelbaarheid van
        \omterm{thm:b3:galois:finitefields}{Frobenius} (vraag 8) voor alle vijf graden. Waar in de
        tafel \emph{zie} je een verschil met $S_5$, waarvan alle
        karakterwaarden gehele rationale getallen zijn?
  \item Leid uit de tafel alleen af dat $A_5$ enkelvoudig is: een
        \omterm{def:b3:groups:normal}{normaaldeler} is een vereniging van \omterm{ex:b3:groups:actions}{conjugatieklassen} die
        $e$ bevat en waarvan de kardinaliteit $60$ deelt ---
        controleer dat geen enkele echte deelsom van $1 + 15 + 20
        + 12 + 12$ die de term $1$ bevat, het getal $60$ deelt.
        Controleer dit tegen het criterium van vraag 11: ga na dat
        geen enkele klasse van $A_5$ priemmachtgrootte $> 1$
        heeft.
  \item (Icosaëdrisch slotstuk) $A_5$ is de rotatiegroep van de
        icosaëder, en de \omterm{def:b3:representations:rep}{representaties} van graad $3$ zijn de twee
        meetkundige \omterm{def:b3:groups:action}{werkingen} op $\R^3$. Ga de spooridentiteit na:
        een rotatie over een hoek $\theta$ heeft spoor $1 +
        2\cos\theta$, en $1 + 2\cos\frac{2\pi}5 =
        \frac{1+\sqrt5}2 = \varphi$. Verklaar zonder enige
        berekening waarom het andere \omterm{def:b3:representations:character}{karakter} van graad $3$ de
        toegevoegde waarde moet dragen: de \omterm{def:b3:galois:galois}{Galoisgroep} van
        $\Q(\sqrt5)/\Q$ werkt op de hele \omterm{def:b3:representations:table}{karaktertafel} (ingang
        voor ingang) en permuteert daarbij de irreducibele
        \omterm{def:b3:representations:character}{karakters}.
\end{enumerate}

\medskip
\noindent\textbf{Deel VI --- Aanvullingen: een centrale grens en
het tensorkwadraat.}
\begin{enumerate}[resume]
  \item (Scherper dan een deelbaarheid) Zij $\chi$ \omterm{def:b3:representations:rep}{irreducibel} van
        graad $n$. Toon aan dat $\abs{\chi(z)} = n$ voor elke $z
        \in Z(G)$ \emph{(lemma van Schur: $\rho(z)$ is een
        scalair, van eindige orde)}, en leid uit $\langle\chi,
        \chi\rangle = 1$ de grens
        \[
        n^2 \leq [G : Z(G)]
        \]
        af. Toon aan dat de niet-abelse groepen van orde $8$ de
        irreducibele graden $1, 1, 1, 1, 2$ hebben \emph{(vijf
        \omterm{ex:b3:groups:actions}{conjugatieklassen}; schrijf $8$ als som van vijf
        kwadraten)} en dat ze de gelijkheid $4 = [G : Z(G)]$
        bereiken; controleer de grens op $A_5$, waarvan het
        \omterm{ex:b3:groups:actions}{centrum} triviaal is.
  \item (Tensorkwadraat van $\chi_2$) Voor $g$ van eindige orde is
        $\rho(g)$ diagonaliseerbaar met eenheidswortels als
        eigenwaarden; leid daaruit de karakterformules
        \[
        \chi_{\operatorname{Sym}^2 V}(g)
        = \frac{\chi(g)^2 + \chi(g^2)}2,
        \qquad
        \chi_{\Lambda^2 V}(g)
        = \frac{\chi(g)^2 - \chi(g^2)}2
        \]
        af. Pas ze toe op $\chi_2$ van $A_5$ \emph{(merk op dat
        $g^2$ de klasse van $c^2$ doorloopt wanneer $g$ die van
        $c$ doorloopt, en omgekeerd)}: toon aan dat
        $\Lambda^2\chi_2 = \chi_2$ en $\operatorname{Sym}^2\chi_2 =
        \mathbf 1 + \chi_5$, zodat
        \[
        \chi_2\otimes\chi_2 = \mathbf 1 + \chi_2 + \chi_5 .
        \]
        Interpreteer $\Lambda^2\chi_2 = \chi_2$ meetkundig via het
        uitwendige product op $\R^3$.
  \item (Slotcontrole van de tafel) Ga numeriek na:
        \omterm{thm:b3:representations:orthogonality}{kolomorthogonaliteit} tussen de twee kolommen van
        $5$-cykels ($\varphi\bar\varphi = -1$), de waarde $5 =
        \abs{Z_{A_5}(c)}$ voor de kolom van $c$ tegen zichzelf, en
        het verdwijnen van het reguliere \omterm{def:b3:representations:character}{karakter} $\sum_i
        n_i\chi_i(g) = 0$ op elk van de vier kolommen van de tafel
        die niet bij de identiteit horen.
\end{enumerate}
\end{problem}
